\chapter{Renormalization}
\label{app:renormalization}

Loop integrals in quantum field theory extend over arbitrarily large
momenta and may therefore diverge.
A consistent treatment of such divergences requires the procedures of
regularization and renormalization.

In this appendix we review the basic framework underlying ultraviolet
divergences, dimensional regularization, and renormalization in
perturbative quantum field theory.
We discuss the structure of counterterms, the role of
one-particle-irreducible (1PI) $n$-point functions and Dyson resummation,
and the definition of renormalization schemes, with particular
emphasis on the $\overline{\mathrm{MS}}$ and on-shell prescriptions.
The discussion follows standard textbook treatments,
in particular Refs.~\cite{PeskinSchroeder1995,WeinbergQFT}.

\section*{Ultraviolet divergences and dimensional regularization}

Consider a generic loop integral of the form
\begin{equation}
I = \int \frac{d^4 p}{(2\pi)^4} \, f(p),
\end{equation}
where $f(p)$ arises from propagators and interaction vertices.
The convergence of the integral is determined by the behavior of
$f(p)$ at large momentum, $p \to \infty$.
If $f(p)$ does not fall off sufficiently fast for large momentum, the integral diverges.

Divergences originating from the high-momentum region are referred to
as \emph{ultraviolet} divergences.
As a simple illustration, an integral behaving as
\begin{equation}
I \sim \int^{\infty} dp \, p^{n}
\end{equation}
diverges for $n \ge -1$.
In four spacetime dimensions, loop integrals typically exhibit either
power-law or logarithmic ultraviolet divergences.

Ultraviolet divergences indicate that perturbative expressions are
formally ill-defined.
To extract meaningful physical results, one introduces a
regularization procedure.
Regularization modifies divergent integrals in a controlled manner so
that they become finite by introducing an additional parameter, for instance
a cutoff scale or a shift in spacetime dimension.
Physical quantities must not depend on this parameter once
renormalization has been performed.

Several regularization methods are commonly used in quantum field theory.
In this thesis we employ dimensional regularization, which preserves
Lorentz invariance and gauge symmetry and is particularly convenient in 
perturbative calculations. The basic idea of dimensional regularization
is to analytically continue the number of spacetime dimensions away from
four. One replaces
\begin{equation}
\int \frac{d^4 p}{(2\pi)^4}
\quad \longrightarrow \quad
\Lambda^{2\epsilon}
\int \frac{d^{d} p}{(2\pi)^{d}},
\qquad
d = 4 - 2\epsilon ,
\end{equation}
where $\epsilon$ is a small parameter and $\Lambda$ is an arbitrary mass
scale introduced to keep the integral dimensionally consistent,
known as the renormalization scale.
The original theory is recovered in the limit $\epsilon \to 0$.

Within dimensional regularization, logarithmic ultraviolet divergences manifest
themselves as poles in the continuation parameter $\epsilon$.
To see how this arises in practice, it is useful to consider a
standard class of Euclidean one-loop integrals.

We define
\begin{equation}
I_\alpha(m^2)
\equiv
\Lambda^{2\epsilon}
\int \frac{d^{d}p}{(2\pi)^{d}}\,
\frac{1}{(p^2+m^2)^\alpha},
\qquad d=4-2\epsilon ,
\label{eq:Ialpha_def}
\end{equation}
where $\alpha$ is a positive parameter.

The evaluation of \eqref{eq:Ialpha_def} proceeds by exploiting
rotational invariance in $d$ dimensions.
Writing the $d$-dimensional momentum measure in spherical coordinates,
one separates the angular and radial integrations according to
\begin{equation}
\int d^d p
=
\Omega_d
\int_0^\infty dp\, p^{d-1},
\end{equation}
where
\begin{equation}
\Omega_d = \frac{2\pi^{d/2}}{\Gamma(d/2)}
\end{equation}
denotes the surface area of the unit sphere in $d$ dimensions.
The remaining radial integral can be evaluated using standard integral
representations of the Gamma function.
One obtains the general result
\cite{PeskinSchroeder1995,WeinbergQFT}
\begin{equation}
I_\alpha(m^2)
=
\frac{\Lambda^{2\epsilon}}{(4\pi)^{d/2}}
\frac{\Gamma\!\left(\alpha-\frac{d}{2}\right)}{\Gamma(\alpha)}
\left(m^2\right)^{\frac{d}{2}-\alpha}.
\label{eq:Ialpha_gamma}
\end{equation}

The Gamma function $\Gamma(z)$ is defined for $\Re(z)>0$ by
\begin{equation}
\Gamma(z)=\int_0^\infty dt\, t^{z-1} e^{-t},
\end{equation}
and extended to other complex values by analytic continuation.
It has simple poles at $z=0,-1,-2,\ldots$.
Logarithmic ultraviolet divergences in dimensional regularization arise
precisely from these poles when the argument of the Gamma function
approaches a non-positive integer.

To make this explicit, consider the case $\alpha=1$.
From the general result \eqref{eq:Ialpha_gamma} we obtain
\begin{equation}
I_1(m^2)
=
\frac{\Lambda^{2\epsilon}}{(4\pi)^{d/2}}
\Gamma\!\left(1-\frac{d}{2}\right)
\left(m^2\right)^{\frac{d}{2}-1}.
\end{equation}

For $d=4-2\epsilon$ the argument of the Gamma function becomes
\begin{equation}
\Gamma\!\left(1-\frac{d}{2}\right)
=
\Gamma(-1+\epsilon).
\end{equation}
The Gamma function has a simple pole at negative integers, and its
expansion around $\epsilon=0$ reads
\begin{equation}
\Gamma(-1+\epsilon)
=
-\frac{1}{\epsilon}
+\mathcal{O}(1).
\end{equation}
Expanding the remaining factors,
\begin{align}
\Lambda^{2\epsilon} &= 1 + 2\epsilon\ln\Lambda + \mathcal{O}(\epsilon^2), \\
(4\pi)^{-d/2} &= (4\pi)^{-2}\left[1+\epsilon\ln(4\pi)+\mathcal{O}(\epsilon^2)\right], \\
(m^2)^{\frac{d}{2}-1} &= m^2\left[1-\epsilon\ln m^2+\mathcal{O}(\epsilon^2)\right],
\end{align}
one finds
\begin{equation}
I_1(m^2)
=
\frac{m^2}{16\pi^2}
\left[
\frac{1}{\epsilon}
+
\ln\!\left(\frac{\Lambda^2}{m^2}\right)
+\text{finite}
\right].
\label{eq:I1_pole}
\end{equation}

The term proportional to $1/\epsilon$ corresponds to the logarithmic
ultraviolet divergence that would appear in four spacetime dimensions.
Dimensional regularization thus converts logarithmic ultraviolet
divergences into simple poles in the analytic continuation parameter
$\epsilon$.

In the main text we instead encounter integrals evaluated in
$d=3-2\epsilon$ spatial dimensions. In this case the argument of the
Gamma function becomes
\begin{equation}
\Gamma\!\left(\alpha-\frac{d}{2}\right)
=
\Gamma\!\left(\alpha-\frac{3}{2}+\epsilon\right).
\end{equation}
Divergences again arise whenever the argument approaches a
non-positive integer. In particular, for $\alpha=2$ one obtains
\begin{equation}
\Gamma\!\left(2-\frac{d}{2}\right)
=
\Gamma\!\left(\frac{1}{2}+\epsilon\right),
\end{equation}
whose expansion generates poles in $\epsilon$ after differentiation
with respect to $m^2$. Thus the ultraviolet divergences appearing in
three spatial dimensions are again represented by poles in the
continuation parameter $\epsilon$.

By contrast, power divergences are not represented in the same way in
dimensional regularization: integrals with no intrinsic scale
(so-called scaleless integrals), such as
\(
\int d^d p\,(p^2)^\beta,
\)
vanish identically. For this reason dimensional regularization
isolates the logarithmic UV divergences—the ones relevant for
renormalization—as poles in $\epsilon$ in a particularly simple
manner.

\paragraph{Remark on $\gamma^5$ in dimensional regularization.}
In theories involving axial currents, special care is required when
extending the Dirac algebra to $d$ dimensions, since the definition
of $\gamma^5$ is intrinsically four-dimensional. In four spacetime
dimensions, $\gamma^5$ is defined as
\begin{equation}
\gamma^5 = i\gamma^0\gamma^1\gamma^2\gamma^3 ,
\end{equation}
but this definition cannot be straightforwardly generalized when the
dimension is analytically continued to $d=4-2\epsilon$.
As discussed in the original work by 't Hooft and Veltman
\cite{tHooftVeltman1972}, ambiguities may arise in loop integrals
involving $\gamma^5$, since not all four-dimensional identities can
be maintained in arbitrary dimensions.

In the present calculation, $\gamma^5$ matrices do appear in the
fermion loop diagrams associated with the pion fields. However, they
always occur in pairs inside Dirac traces, such that they can be
combined using $(\gamma^5)^2 = 1$ before performing the dimensional
continuation. As a result, no traces involving an odd number of
$\gamma^5$ matrices are encountered, and the potentially ambiguous
structures do not arise.

We therefore assume that $\gamma^5$ can be treated as anticommuting
with $\gamma^\mu$ and satisfying $(\gamma^5)^2=1$ for the purpose of
these calculations. This is sufficient at the present order, since
the problematic aspects of $\gamma^5$ in dimensional regularization
do not enter.

The remaining Dirac identities used in the calculation, such as
$\mathrm{Tr}\,\gamma^\mu = 0$ and
$\mathrm{Tr}\,(\gamma^\mu\gamma^\nu)=4g^{\mu\nu}$, are unaffected by
the dimensional continuation.

\subsection*{Useful integrals}
\label{app:loop_integrals}

In the main text, the renormalization of the Quark--Meson model
involves a small set of standard one-loop integrals.
These can be expressed in terms of the general result
Eq.~\eqref{eq:Ialpha_gamma}.

The tadpole integral is defined as
\begin{equation}
A(M_q^2)
=
\Lambda^{4-d}
\int \frac{d^d k}{(2\pi)^d}\,
\frac{1}{k^2 - M_q^2 + i\epsilon},
\end{equation}
where $d=4-2\epsilon$.
Using Eq.~\eqref{eq:Ialpha_gamma} with $\alpha=1$, one obtains
\begin{equation}
A(M_q^2)
=
\frac{i M_q^2}{(4\pi)^2}
\left[
\frac{1}{\epsilon}
+
\ln\!\left(\frac{\Lambda^2}{M_q^2}\right)
+ 1
\right].
\end{equation}

The two-propagator integral is defined by
\begin{equation}
B(p^2)
=
\Lambda^{4-d}
\int \frac{d^d k}{(2\pi)^d}\,
\frac{1}{\bigl[(k+p)^2 - M_q^2 + i\epsilon\bigr]
\bigl[k^2 - M_q^2 + i\epsilon\bigr]}.
\end{equation}
Using standard Feynman parameterization and dimensional
regularization, this integral can be written as
\begin{equation}
B(p^2)
=
\frac{i}{(4\pi)^2}
\left[
\frac{1}{\epsilon}
+
F(p^2)
\right],
\end{equation}
where the finite function $F(p^2)$ is given by
\begin{equation}
F(p^2)
=
2
-
2 \sqrt{\frac{4M_q^2}{p^2} - 1}
\arctan\!\left(
\frac{1}{\sqrt{\frac{4M_q^2}{p^2} - 1}}
\right),
\end{equation}
for $p^2 < 4M_q^2$.

We also define
\begin{equation}
B'(p^2)
=
\frac{dB(p^2)}{dp^2}
=
\frac{i}{(4\pi)^2} F'(p^2),
\end{equation}
which appears in the wavefunction renormalization.

\section*{Renormalization and counterterms}

Regularization renders loop integrals finite by introducing an auxiliary
parameter, such as the continuation parameter $\epsilon$ in dimensional
regularization.
However, the regulated expressions still contain poles as
$\epsilon \to 0$.
Renormalization is the procedure by which these divergences are absorbed
into redefinitions of the parameters and fields of the theory, so that
physical quantities remain finite in the limit
$\epsilon \to 0$.
For general discussions see, for example,
Refs.~\cite{Dyson1949,PeskinSchroeder1995,WeinbergQFT}.

The modern understanding of renormalization emerged from the analysis
of ultraviolet divergences in quantum electrodynamics in the late 1940s,
notably in the works of Dyson, Schwinger, Tomonaga, and Feynman
\cite{Dyson1949,Schwinger1948,Tomonaga1946,Feynman1949}.
Dyson demonstrated the equivalence of the various formulations of quantum
electrodynamics and clarified the perturbative structure of renormalization
\cite{Dyson1949}.
The systematic interpretation of renormalization as a redefinition of
fields and parameters was subsequently developed in the framework of
relativistic quantum field theory.

In a renormalizable theory, the Lagrangian is written in terms of
\emph{bare} fields and \emph{bare} parameters.
For a scalar field $\phi$, one introduces a wavefunction renormalization
factor $Z_\phi$ such that
\begin{equation}
\phi_0 = Z_\phi^{1/2}\phi ,
\end{equation}
where $\phi_0$ denotes the bare field and $\phi$ the renormalized field.
Similarly, bare parameters are expressed in terms of renormalized
parameters and counterterms,
\begin{equation}
m_0^2 = m^2 + \delta m^2,
\qquad
g_0 = g + \delta g .
\end{equation}
The counterterms $\delta m^2$, $\delta g$, and
$\delta Z_\phi \equiv Z_\phi - 1$
are introduced precisely so that ultraviolet divergences can be absorbed
into these redefinitions.

Substituting these relations into the bare Lagrangian and expanding to
first order in the counterterms separates the theory into two parts.
The first part has the same functional form as the original Lagrangian,
but is expressed entirely in terms of renormalized quantities.
The second part consists of additional interaction terms proportional
to $\delta Z_\phi$, $\delta m^2$, $\delta g$, etc.
These additional terms are referred to as \emph{counterterm vertices}.
They are treated perturbatively alongside the original interaction
vertices.

The structure of renormalization becomes particularly transparent when
considering Green’s functions.
Loop diagrams generate ultraviolet divergences in $n$-point functions.
The counterterm vertices produce additional contributions with exactly
the same Lorentz and momentum structure as the divergent parts of the
loop corrections.
By choosing the counterterms appropriately, the sum of loop diagrams and
counterterm diagrams becomes finite as $\epsilon \to 0$.
This cancellation is implemented order by order in perturbation theory.

The possibility of absorbing all ultraviolet divergences into a finite
number of redefinitions of fields and parameters is the defining property
of a renormalizable theory.
Power-counting arguments establish which interactions satisfy this
criterion in four spacetime dimensions
\cite{WeinbergQFT}.
More generally, renormalization can be understood as a consequence of
the short-distance structure of quantum field theory and its behavior
under changes of scale, as emphasized in the renormalization group
approach \cite{GellMannLow1954,Wilson1971}.

\section*{Perturbation theory, 1PI functions, and Dyson resummation}

In perturbative quantum field theory, Green's functions are expanded
in powers of the interaction coupling.
Each term in this expansion can be represented graphically by a
Feynman diagram, following the rules derived from the interaction
Lagrangian \cite{Feynman1949,PeskinSchroeder1995,WeinbergQFT}.
Internal lines correspond to free propagators, vertices correspond
to interaction terms, and each independent loop momentum is integrated
over.

An important class of diagrams is formed by the
\emph{one-particle-irreducible} (1PI) diagrams.
A diagram is said to be 1PI if it cannot be disconnected by cutting
a single internal propagator line.
The 1PI $n$-point functions, usually denoted
$\Gamma^{(n)}$, play a central role in renormalization theory
\cite{WeinbergQFT}.

For the two-point function, the sum of all 1PI diagrams defines
the self-energy $\Sigma(p^2)$.
Diagrammatically, the full propagator can be expressed as a series of
repeated insertions of the one-particle-irreducible two-point function,
as illustrated in Fig.~\ref{fig:dyson_series}. The circle labelled
$1\mathrm{PI}$ denotes the sum of all one-particle-irreducible
two-point diagrams. Starting from the free propagator
\begin{equation}
G_0(p) = \frac{i}{p^2 - m_0^2 + i\epsilon},
\end{equation}
successive insertions of the self-energy generate a geometric series
\begin{equation}
G(p)
=
G_0(p) + G_0(p)\Sigma(p^2)G_0(p)
+ G_0(p)\Sigma(p^2)G_0(p)\Sigma(p^2)G_0(p) + \cdots .
\end{equation}
Resumming this series yields Dyson's equation
\cite{Dyson1949},
\begin{equation}
G(p)
=
\frac{i}{p^2 - m_0^2 - \Sigma(p^2) + i\epsilon}.
\label{eq:Dyson_appendix}
\end{equation}

\begin{figure}[htbp]
\centering
\begin{tikzpicture}
\begin{feynman}
\vertex (L1) at (-6.8,0);
\vertex (L2) at (-4.4,0);
\vertex[draw, circle, fill=gray!25, minimum size=12mm] (Full) at (-5.6,0) {};
\diagram*{
  (L1) -- (Full) -- (L2),
};
\node at (-3.8,0) {\(=\)};
\vertex (a1) at (-3.2,0);
\vertex (a2) at (-2.0,0);
\diagram*{
  (a1) -- (a2),
};
\node at (-1.4,0) {\(+\)};
\vertex (b1) at (-0.8,0);
\vertex (b2) at (1.2,0);
\vertex[draw, circle, minimum size=10mm, inner sep=1pt] (Sigma1) at (0.2,0) {\(1\mathrm{PI}\)};
\diagram*{
  (b1) -- (Sigma1) -- (b2),
};
\node at (1.8,0) {\(+\)};
\vertex (c1) at (2.4,0);
\vertex (c2) at (5.8,0);
\vertex[draw, circle, minimum size=10mm, inner sep=1pt] (Sigma2) at (3.4,0) {\(1\mathrm{PI}\)};
\vertex[draw, circle, minimum size=10mm, inner sep=1pt] (Sigma3) at (4.8,0) {\(1\mathrm{PI}\)};
\diagram*{
  (c1) -- (Sigma2) -- (Sigma3) -- (c2),
};
\node at (6.6,0) {\(+\ \cdots\)};
\end{feynman}
\end{tikzpicture}
\caption{Diagrammatic Dyson series for the two-point function. The
circle labelled $1\mathrm{PI}$ represents the sum of all
one-particle-irreducible two-point diagrams, i.e.\ the self-energy
$\Sigma(p^2)$.}
\label{fig:dyson_series}
\end{figure}

Equation~\eqref{eq:Dyson_appendix} shows that the self-energy acts as
a momentum-dependent correction to the inverse propagator.
Loop effects therefore modify both the position of the propagator pole
and its residue.

It is convenient to separate the bare mass parameter $m_0^2$ into a
renormalized mass and a counterterm,
\begin{equation}
m_0^2 = m^2 + \delta m^2 ,
\end{equation}
and to include the wavefunction renormalization factor introduced in
the previous subsection.
The renormalized self-energy $\Sigma^{\mathrm{ren}}(p^2)$ is then defined
as the sum of loop diagrams and counterterm contributions.
Renormalization conditions determine the counterterms so that
$G(p)$ has the desired physical properties.

\section*{Renormalization schemes: MS, $\overline{\mathbf{MS}}$, and on-shell}

The introduction of counterterms removes the ultraviolet divergences
appearing in loop diagrams.
However, the divergent parts do not uniquely determine the finite parts
of the counterterms.
The prescription used to fix these finite contributions defines a
\emph{renormalization scheme}
\cite{PeskinSchroeder1995,WeinbergQFT}.

Different renormalization schemes correspond to different choices of
finite counterterms.
While intermediate quantities such as Green's functions and
renormalized parameters depend on this choice, physical observables
are independent of the renormalization scheme when computed exactly.

\paragraph{Minimal subtraction (MS).}

In the minimal subtraction scheme, one removes only the pole terms
in $\epsilon$ that arise in dimensional regularization.
If a loop integral produces a divergence of the form
\begin{equation}
\frac{1}{\epsilon} + \text{finite},
\end{equation}
the counterterm is chosen to cancel only the $1/\epsilon$ part.
No additional finite terms are subtracted.

In practice, it is common to use the modified minimal subtraction
scheme, denoted $\overline{\mathrm{MS}}$.
In this scheme one subtracts, together with the pole,
the universal constants that arise from the expansion of Gamma
functions in dimensional regularization,
\begin{equation}
\frac{1}{\epsilon}
\;\longrightarrow\;
\frac{1}{\epsilon}
- \gamma_E
+ \ln(4\pi),
\end{equation}
where $\gamma_E$ is the Euler–Mascheroni constant.
The $\overline{\mathrm{MS}}$ prescription simplifies many perturbative
expressions and is therefore widely used in perturbative quantum field
theory, including quantum chromodynamics.

In minimal subtraction schemes, renormalized parameters depend
explicitly on the renormalization scale $\Lambda$ introduced in
dimensional regularization.
This scale dependence is governed by renormalization group equations,
first derived in the context of quantum electrodynamics by
Gell-Mann and Low \cite{GellMannLow1954} and later developed in
a general framework by Wilson \cite{Wilson1971}.

\paragraph{On-shell scheme.}

In contrast to MS-type schemes, which subtract divergences according
to an algebraic rule, the on-shell scheme fixes counterterms by
imposing physical renormalization conditions directly on Green’s
functions \cite{Dyson1949,PeskinSchroeder1995,WeinbergQFT}.

For a scalar field, the full propagator has the form
\begin{equation}
G(p^2)
=
\frac{i}{p^2 - m^2 - \Sigma^{\mathrm{ren}}(p^2) + i\epsilon}.
\end{equation}
In the on-shell scheme, the renormalized mass parameter $m$ is defined
as the physical pole mass.
This requirement implies the condition
\begin{equation}
\Sigma^{\mathrm{ren}}(m^2) = 0.
\label{eq:OS_mass_scheme}
\end{equation}

A second condition fixes the residue of the propagator at the pole.
Expanding the denominator around $p^2 = m^2$,
one finds that the residue equals
\(
\left[1 - \Sigma^{\prime\,\mathrm{ren}}(m^2)\right]^{-1}.
\)
The on-shell scheme imposes
\begin{equation}
\Sigma^{\prime\,\mathrm{ren}}(m^2) = 0,
\label{eq:OS_residue_scheme}
\end{equation}
so that near the pole
\begin{equation}
G(p^2) \sim \frac{i}{p^2 - m^2}.
\end{equation}
This condition determines the wavefunction renormalization constant
and ensures canonical normalization of one-particle states.
